% Generated draft: requires mathematical and source review.
\documentclass[11pt]{article}
\usepackage[margin=25mm]{geometry}
\usepackage{fontspec}
\setmainfont{Latin Modern Roman}
\usepackage{amsmath,amssymb,mathtools,unicode-math}
\setmathfont{Latin Modern Math}
\usepackage{enumitem,needspace,xurl,hyperref}
\setlength{\parindent}{0pt}
\setlength{\parskip}{5pt}
\setlength{\emergencystretch}{3em}
\newcommand{\N}{\mathbb{N}}
\newcommand{\Z}{\mathbb{Z}}
\newcommand{\Q}{\mathbb{Q}}
\newcommand{\R}{\mathbb{R}}
\newcommand{\C}{\mathbb{C}}
\newcommand{\F}{\mathbb{F}}
\newcommand{\A}{\mathbb{A}}
\newcommand{\PP}{\mathbb{P}}
\newcommand{\E}{\mathbb{E}}
\newcommand{\eps}{\varepsilon}
\newcommand{\dd}{\,\mathrm{d}}
\begin{document}
\section*{UnsolvedMath ID: 11000190}
\textbf{Source problem code:} catalogue entry
\subsection{Definitions and mathematical statement}
\textbf{Conjecture 2.3, as extracted.} If $r\ge3$, the action of $\operatorname{Out}(\pi)$ on $\operatorname{Hom}(\pi,G)$ is ergodic. The excerpt does not restate what $r$, $\pi$, $G$, or the invariant measure are, and its literal action notation needs clarification.

The customary free-group interpretation is $\pi=F_r$, a free group of rank $r$, with $G$ a compact connected Lie group. A basis identifies $\operatorname{Hom}(F_r,G)$ with $G^r$, equipped with product normalized Haar measure. Automorphisms of $F_r$ act by precomposition. Inner automorphisms of $F_r$ generally change a representation by conjugation in $G$, so the natural outer-automorphism action is on the quotient
\[
 X_r(G)=\operatorname{Hom}(F_r,G)/G,
\]
with the pushforward Haar probability measure $\mu$, not on unquotiented $\operatorname{Hom}(F_r,G)$ in general. Under this interpretation the precise proposed assertion is
\[
 \forall r\ge3,\qquad
 A\subset X_r(G)\text{ measurable and }\operatorname{Out}(F_r)\text{-invariant}
 \ \Longrightarrow\ \mu(A)\in\{0,1\}.
\]
The allowed class of compact groups must be confirmed against the surrounding source; this is not a conjecture for arbitrary groups with an unspecified measure. An alternative literal representation-space assertion would use $\operatorname{Aut}(F_r)$, and would be a different formulation requiring explicit authorization from the source text. Neither the quotient nor the acting group has been silently corrected.

The source reports that calculations in [43] were used to prove the conjecture in [47] when every simple factor of $G$ is locally isomorphic to $\mathrm{SU}(2)$. That restricted-case claim is retained with its source attribution and does not establish the general compact-group assertion or resolve the notation defect.
\subsection{Short English statement}
The likely conjecture is ergodicity of the outer-automorphism action on the compact-group character space of a free group of rank at least three. The literal excerpt omits the quotient and measure, which are essential to a well-defined statement.
\subsection{Sources}
\begin{itemize}
\item Local UnsolvedMath registry, permanent ID 11000190, source code catalogue entry; no individual proposer or proposal year supplied.
\item Farb, edited volume \emph{Problems on Mapping Class Groups and Related Topics}, chapter ``Mapping class group dynamics on surface group representations,'' Conjecture 2.3, PDF page 217; \url{https://www.math.uchicago.edu/~farb/papers/mcgbook.pdf}. The registry records PDF-font-structure extraction and access on 2026-07-29; its inventory distinguishes 329 labeled headings by source page.
\item Source references [43] and [47], respectively for calculations and the locally-$\mathrm{SU}(2)$ simple-factor case. Full bibliography entries and exact hypotheses are absent.
\item Adjacent local ID 11000189 mentions compact nonabelian groups, but does not reproduce the full free-group setup of this conjecture.
\end{itemize}
\end{document}
